%% 第二章 相关工作综述

\chapter{相关工作综述}
\chaptermark{相关工作综述}

\section{引言}

长期机器人经验问答的核心在于，如何对机器人在运行过程中产生的大量多模态经历进行有效组织、检索与维护。与传统短时任务不同，长期部署场景中的机器人需要面对跨时间跨度的历史信息管理问题，包括事件表示、层级摘要、语义检索、记忆压缩以及错误修正等多个方面。因此，构建具备长期记忆能力的情景记忆系统，已经成为当前具身智能与服务机器人研究中的重要方向。

机器人长期情景记忆涉及记忆表示、经验组织、信息检索、长期压缩与动态修正等多个问题。近年来，研究者围绕自然语言记忆表示、层级式经验组织、检索增强生成以及长期视频问答等方向开展了大量研究工作，为机器人长期经验问答提供了重要基础。

本章围绕本文所涉及的关键技术和相关工作，对机器人情景记忆、层级记忆组织、图增强检索、长期记忆管理以及视频理解等相关研究进行综述，并分析现有工作的特点与局限，为后续章节的方法设计提供技术背景。

\section{机器人情景记忆}

情景记忆（Episodic Memory）的概念来源于认知科学，用于描述智能体对过去事件及其上下文信息的记忆能力。近年来，情景记忆已被广泛应用于智能穿戴设备、智能空间以及机器人系统等领域。在机器人场景中，情景记忆通常用于记录机器人在长期运行过程中产生的感知、动作与交互历史，从而支持后续的问答、推理与任务规划。

现有机器人情景记忆方法大致可分为三类。

第一类方法基于隐向量表示。该类方法将机器人的经历编码为固定维度的潜在向量。这类方法通常使用循环神经网络（RNN）或变分自编码器（VAE）对时序化的感知数据进行压缩编码，在检索时通过向量相似度匹配相关经历\cite{elman1990rnn,kingma2014vae,hochreiter1997lstm,bahdanau2015attention}。例如，文献\cite{hochreiter1997lstm}使用长短期记忆网络（LSTM）将机器人操作序列编码为情景向量，在后续任务中通过解码器重构相关经历；文献\cite{bahdanau2015attention}进一步引入注意力机制，使编码器能够选择性关注经历中的关键帧。这类方法的优势在于检索效率高，且端到端训练使得编码器和下游任务可以联合优化。然而，隐向量方法的局限也十分显著：隐向量的语义不可直接阅读，用户和开发者无法检查"机器人记住了什么"；编码器针对特定任务和场景训练，当机器人遇到训练数据之外的新物体或新任务时，编码质量显著下降；固定维度的潜在向量在压缩信息时不可避免地丢失细节，一个 512 维的向量无法完整保留事件中涉及的时间、地点、物体属性、动作及状态变化等全部信息维度。

\begin{figure}[htbp]
    \centering
    \includegraphics[width=0.85\textwidth]{图2-1基于隐向量表示的机器人情景记忆方法示意图.png}
    \caption{基于隐向量表示的机器人情景记忆方法示意图}
    \label{fig:chapter2-latent}
\end{figure}

第二类方法采用显式符号化存储。该方法将经历表示为结构化的符号记录（如键值对、关系图谱或一阶逻辑命题），并通过基于规则的检索接口进行查询。代表性工作如 KnowRob\cite{tenorth2013knowrob} 构建了基于本体的机器人知识库，通过逻辑推理引擎支持对物体位置、任务状态和事件时序的查询；RoboEarth\cite{waibel2011roboearth} 则提出了云端的机器人经验共享平台，不同机器人可以将自身经历以结构化语义地图的形式上传和检索。符号化方法的核心优势在于精确性和可解释性，检索结果是确定性的，推理链可追踪。但其最突出的瓶颈在于可扩展性：规则库和本体的构建高度依赖人工设计，当环境中的物体类别从几十扩展到数千、任务类型从预定义扩展到开放域时，手工维护知识体系的工作量呈指数级增长。此外，符号化方法难以处理自然语言中的模糊性和多义性，同一问题可能有多种语义等价的表述方式，而规则式检索无法灵活应对。

基于自然语言表示的大模型驱动方法，是近年来大语言模型的涌现能力开辟的第三条路径：该类方法直接将机器人经历转化为文本描述，利用预训练大模型的语义理解和生成能力实现灵活的记忆检索和问答\cite{liu2023reflect}。例如，REFLECT 等工作使用自然语言对机器人经验进行组织，并通过语言模型实现长期任务回顾与行为分析。相比传统符号方法，自然语言表示兼具一定的可解释性与语义泛化能力，更适合开放环境中的长期交互任务。

在此基础上，H-EMV（Hierarchical Episodic Memory for Long-Term Robot Experience Question Answering）进一步提出层级化情景记忆结构，将机器人长期经历组织为多层树状结构，并利用大语言模型以交互式方式在记忆树中逐层检索相关信息。该方法能够在较低 token 开销下完成长时间跨度的机器人经验问答，在 TEACh、Ego4D 和 ARMAR-7 等数据集上取得了较好的效果。

总体来看，现有机器人情景记忆方法已经能够实现长期经历的组织与问答，但在跨事件关联、长期记忆维护以及错误修正等方面仍有进一步提升空间。

\begin{table}[htbp]
\centering
\caption{机器人情景记忆三阶段方法对比}
\label{tab:2-1}
\begin{tabular}{|p{2.0cm}|p{2.5cm}|p{2.5cm}|p{1.2cm}|p{2.0cm}|p{1.2cm}|}
\hline
方法类型 & 记忆表示 & 检索方式 & 可解释性 & 泛化性 & 可维护性 \\
\hline
隐向量方法 & 固定维度潜在向量 & 向量相似度 & 低 & 受限于训练数据 & 无 \\
\hline
符号存储方法 & 结构化键值/图谱 & 规则推理 & 高 & 受限于手工本体 & 无 \\
\hline
大模型驱动方法 & 自然语言文本 & LLM Agent 交互式 & 高 & 利用预训练模型 & 弱 \\
\hline
\end{tabular}
\end{table}

\section{层级情景记忆框架}

本文所采用的 GRAF-Mem 框架建立在层级情景记忆结构之上，因此有必要对其基本组织方式进行介绍。

H-EMV 使用层级记忆树组织机器人长期经历。该结构从底层原始感知数据逐层抽象形成高层语义摘要。其中，底层节点保存图像、音频、场景图等多模态信息，中间层记录事件级与目标级摘要，高层节点则对更长时间范围内的经历进行递归概括。通过这种层级化组织方式，系统能够在保留局部细节的同时，对长期历史进行压缩表示。

\begin{figure}[htbp]
    \centering
    \includegraphics[width=\textwidth]{图2-2 层级情景记忆树表示.png}
    \caption{层级情景记忆树表示}
    \label{fig:chapter2-hierarchical}
\end{figure}

与传统一次性检索方式不同，H-EMV 将大语言模型作为交互式 Agent 使用。系统首先在高层摘要中定位相关时间段，再逐步展开子节点以获取更细粒度的信息，从而形成类似"逐层回忆"的检索过程。这一循环模拟了人类在回忆时的逐层聚焦过程——先回想宏观的时间段，再逐步定位到具体事件。这种方式能够减少无关历史信息带来的上下文开销，并提高长历史场景下的问答效率。

此外，H-EMV 的层级结构具有较好的可解释性。模型在检索过程中形成的树导航路径，可以反映系统如何逐步定位相关经历，从而为最终答案提供一定的推理依据。

然而，现有层级情景记忆框架主要依赖树结构中的父子关系以及文本相似度进行检索，对于跨事件、跨时间段的关联信息利用仍较有限。同时，已有方法大多采用追加式存储机制，对长期运行中的记忆压缩、遗忘与修正问题关注相对较少。

\section{记忆检索与图增强方法}

基于向量相似度的记忆检索是当前大模型应用中最主流的检索范式，代表性框架为检索增强生成（Retrieval-Augmented Generation, RAG）\cite{lewis2020rag}。其核心流程为：将记忆条目切分为文本块，使用嵌入模型（如 Sentence-BERT\cite{reimers2019sbert}）将各文本块编码为固定维度的稠密向量，构建向量索引（如 FAISS\cite{johnson2019faiss}）；检索时将查询编码为同维度向量，通过近似最近邻搜索返回与查询最相似的 top-k 个文本块，送入大语言模型生成最终答案。为了进一步提升检索效果，近年来研究者提出了多种改进方法，如 HyDE\cite{gao2023hyde} 在检索前使用大模型生成假设性答案以缩短查询与文档的语义距离，Re-ranking\cite{nogueira2019reranking} 在初步检索后使用更强的交叉编码器对候选结果进行精排，以提高最终召回质量。

在长期机器人情景记忆场景下，仅依赖文本相似度的检索方式仍存在一定局限。一方面，同一事件可能具有不同语言表达形式，仅依赖词汇层面的相似度容易出现召回不足；另一方面，传统向量检索通常将各记忆片段视为相互独立的文本块，缺少事件之间的结构化关联信息。

然而，在机器人长期情景记忆场景下，纯向量检索面临两个核心挑战：（1）措辞鸿沟：同一事件或物体可能以完全不同的措辞被描述，例如用户问"那个蓝色的杯子"，而记忆中的相关事件可能仅描述为 "picked up a mug from the sink"，纯文本相似度无法跨越这种语义等价但词汇层面不相似的鸿沟；（2）横向关联缺失：向量搜索是 item-to-query 的独立匹配，返回的是离散的文本块列表，不携带任何块间关系信息，缺少事件之间的结构化关联信息，无法回答"这个杯子后来出现在哪些任务中"等需要沿时间线追踪同一物体跨任务出现记录的查询。

\begin{figure}[htbp]
    \centering
    \includegraphics[width=\textwidth]{图2-3纯向量检索和图增强检索对比图.png}
    \caption{纯向量检索和图增强检索对比图}
    \label{fig:chapter2-retrieval}
\end{figure}

为弥补纯向量检索在结构化关联捕捉上的不足，近年来出现了多种将图结构引入检索流程的方法。知识图谱增强的 RAG（Graph RAG）\cite{edge2024graphrag} 在检索前利用知识图谱中的实体关系对查询进行扩展，例如从 "blue cup" 扩展至 "mug, cup, vessel" 等上位词/同义词，以提升召回率；层次化索引方法（如 RAPTOR\cite{sarthi2024raptor}）将文档组织为聚类树，在不同层级上分别进行向量检索和文本生成，使检索结果既能覆盖宏观主题也能命中微观细节；在视频理解领域，VideoTree\cite{wang2024videotree} 和 VideoAgent\cite{fan2024videoagent} 等工作也采用类似的"结构化索引 + LLM Agent 导航"范式，将长视频的帧序列组织为树或图，让 Agent 通过自适应展开定位关键帧。

此外，一些研究开始将大语言模型作为交互式 Agent 使用。例如，ReMEmbR 利用语言模型在向量记忆库中主动检索机器人经历；VideoAgent 等方法则允许语言模型通过多轮交互动态获取视频内容信息。这类方法不再依赖固定检索流程，而是通过语言模型自主决定后续的信息访问路径。

总体来看，图增强与层级化检索方法为长期机器人记忆系统提供了新的思路，但现有方法大多侧重于离线索引组织，对机器人长期经历中的时间关联、物体关联与任务关联利用仍不够充分。

\section{长期记忆管理与修正}

随着机器人长期运行时间不断增长，如何在有限存储与上下文资源下管理大量历史经历，逐渐成为长期记忆系统的重要问题。

传统遗忘机制主要来源于认知科学中的遗忘理论以及计算机系统中的缓存管理策略。例如，基于时间衰减的遗忘模型、LRU 缓存淘汰策略以及滑动窗口方法等，均通过时间或访问频率决定信息保留程度。这类方法实现简单，但通常难以区分不同事件的重要性。

经典遗忘模型艾宾浩斯遗忘曲线\cite{ebbinghaus1885memory}指出记忆保持率随时间呈指数衰减。在信息检索和数据库管理系统中，这一原则被广泛应用于缓存置换策略（如 LRU 及其变体）和时间窗口策略（如滑动窗口、衰减窗口），前者优先淘汰最近最少使用的条目，后者通过设置固定的保留时长阈值自动删除超过时限的数据。然而，这些经典策略均通过时间或访问频率决定信息保留程度，在面对机器人情景记忆时存在明显不足：时间不是决定记忆价值的唯一因素，三周前的一次失败任务可能比三天前的一次成功例行操作对后续问答更重要；固定时间窗无法区分不同信息类型，图像等大对象应比文本更激进地压缩，但文本中的关键语义信息需要保留。

随着大语言模型上下文窗口的不断扩展，如何在有限上下文内有效管理长对话/长文档中的记忆信息成为研究热点。MemGPT\cite{packer2023memgpt} 提出了一种受操作系统虚拟内存管理启发的分层记忆架构——将 LLM 的上下文视为"主内存"，将外部存储视为"虚拟内存"，通过类似页面置换的机制在对话过程中自动将历史信息在主内存和虚拟内存之间换入换出。MemWalker\cite{chen2023memwalker} 将长文本组织为摘要树，在问答时让 LLM Agent 从根节点开始逐层导航至相关叶节点，与 GRAF-Mem 的树导航机制有相似之处，但其树结构是静态的且不支持压缩。LongMem\cite{wang2023longmem} 则在 Transformer 解码器之上增加了一个侧边记忆模块，通过注意力机制实现对历史信息的软寻址和选择性遗忘。

这些工作为长期记忆管理提供了重要参考，但也存在关键差异：它们的管理对象是大语言模型上下文中的文本序列，而本文的管理对象是包含结构化的多模态层级数据的记忆树——压缩决策需要感知节点的数据类型和层级位置；此外，这些方法的记忆管理多服务于单次长对话中的"当前任务"而非跨越数周到数月的"长期知识"，因此缺少对新近性和重要性的联合评估机制。

除记忆压缩外，长期运行中的错误修正同样是重要问题。当前相关研究主要包括两类方向。

第一类方法通过用户反馈对模型参数进行更新，在线强化学习\cite{singh2000rl} 将用户的正/负反馈作为奖励信号，通过策略梯度方法更新对话策略的参数；Active Retrieval\cite{jiang2023activeretrieval} 等工作利用用户反馈信号调整检索查询的生成策略。这类方法的优势在于可以通过持续交互不断优化系统表现，但每次修正都需要经过模型推理和参数更新，计算成本高，且参数更新可能带来全局性的副作用。

第二类方法为知识编辑（Knowledge Editing）。代表性方法如 ROME\cite{meng2022rome} 和 MEMIT\cite{meng2023memit} 通过对 Transformer 前馈网络中的特定神经元（通常位于 MLP 层的中间层，被认为存储了事实性知识）进行精确的权重修改，实现对特定事实条目的定位和修正。这类方法的优势在于修改的精确性和局部性：理想情况下只改变目标知识而不影响其他知识。但在实际应用中，知识编辑仍面临泛化性不足（修改后的知识未必能在不同的提问方式下稳定输出正确回答）和副作用评估困难（难以穷举被修改神经元所影响的所有知识条目）等挑战。

上述修正方法均在模型参数空间或检索策略空间中操作，不同的是，机器人长期情景记忆中的错误往往来源于记忆数据本身，例如场景识别偏差或事件摘要错误。因此，相比直接修改模型参数，对记忆条目进行局部修正更加适合长期机器人记忆场景。

\section{视频理解与长视频问答}

视频理解，尤其是视频问答（Video Question Answering，VideoQA）任务，与机器人情景记忆具有较强相关性。二者都需要在长时间数据流中定位关键信息，并基于历史内容完成推理与回答。

近年来，长视频理解和第一人称视频理解成为 VideoQA 的重要研究方向。Ego4D 是当前具有代表性的第一人称视频数据集，记录了大量真实日常活动场景。QAEgo4D、EgoSchema 等基准进一步推动了长时序视频问答研究的发展。

现有 VideoQA 方法大致可以分为两类。一类为端到端训练方法，通过视觉编码器与语言模型联合建模视频内容；另一类为训练自由（training-free）方法，利用视频字幕、视觉标签或外部工具将视频转换为文本描述，再借助大语言模型完成推理与问答。

例如，VideoTree 使用树状结构组织长视频内容，并通过自顶向下的方式逐步定位关键帧；一些工作则将大语言模型作为交互式 Agent，通过多轮调用视觉工具分析视频内容。这类方法与机器人长期情景记忆中的层级检索思路具有一定相似性。

然而，长视频问答中的"长时序"与机器人长期情景记忆中的"长期记忆"仍存在明显差异。现有长视频理解任务中的视频长度通常为数十分钟至数小时，其核心挑战主要在于连续视频流中的关键片段定位与上下文压缩。而机器人长期情景记忆往往需要处理跨越数天、数周甚至数月的离散经历数据，不同事件之间可能存在较大的时间间隔，同时还伴随任务切换、环境变化以及跨场景关联等问题。因此，机器人长期记忆不仅需要解决长上下文建模问题，还需要进一步考虑长期存储、跨时间检索以及动态记忆维护等更复杂的挑战。

此外，机器人情景记忆不仅包含视觉信息，还涉及动作执行、任务状态、环境交互以及真实时间信息等多模态内容。因此，相比一般的视频问答任务，长期机器人经验问答在数据组织方式和长期时间建模方面具有更高复杂度。

\section{本章小结}

本章围绕机器人长期情景记忆相关技术和工作进行了综述。首先介绍了机器人情景记忆的发展过程，分析了隐向量表示、符号化存储以及自然语言记忆表示等不同技术路线的特点；随后对 H-EMV 层级情景记忆框架进行了介绍，并总结了层级化经验组织与交互式检索的基本思想。

在此基础上，本章进一步回顾了图增强检索、长期记忆管理、知识修正以及长视频问答等相关研究方向。总体来看，现有工作已经为长期机器人经验问答提供了重要基础，但在跨事件关联、长期记忆维护以及动态修正等方面仍存在进一步研究空间。

后续章节将在上述研究基础上，分别介绍本文提出的事件关联图机制、效用引导记忆管理方法以及反馈驱动记忆修正机制。
